\section{从一个问题开始}
\subsection{如何给出常数预测？}
\begin{redbox}{本节目标}
理解“所有样本预测同一个数”的含义，计算训练目标的均值。
本节是展示模板的自编短例，不对应原教材习题。
\end{redbox}
\begin{examplebox}{业务例子}
三条虚构的历史记录，其目标值为18、22、26。
如果暂时不使用个体特征，可以先把训练目标均值作为基准预测。
\end{examplebox}
均值为 $\bar y=(18+22+26)/3=22$。这个数来自训练记录。
\CodeBlock{111-01}{计算常数预测值}
\RunOutput{111-01}
\Thought{为什么不能拿测试目标的均值当作预测值？}
请先自己回答，再阅读下面的参考答案。
\begin{answerbox}
这样会让测试目标参与构造预测规则。用于最终评价的数据应保留到方案固定以后；
本例的常数应从训练目标中计算。
\end{answerbox}
